\subsection{Convolution of a measure and a function}

Let X be a locally compact space on which a locally compact group G operates on the left continuously. Let \(\beta\) be a positive measure on X, quasi-invariant under G. Let \(\chi\) be a function \textgreater0 on \(G \times X\) , measurable for every measure on \(\mathbf{G} \times \mathbf{X}\) , and such that, for every \(s \in \mathbf{G}\) , \(\chi(s^{-1}, \cdot)\) is a density of \(\pmb{\gamma}(s)\beta\) with respect to \(\beta\) :

\[
\pmb {\gamma} (s) \beta = \chi (s ^ {- 1}, \cdot) \cdot \beta ,\tag{1}
\]

which, with the conventions of Ch. VII, §1, No.~1, may be written:

\[
d \beta (s x) = \chi (s, x) d \beta (x).\tag{1'}
\]

These data will remain fixed in Nos. 1, 2, 3 (an exception being made in Remark 2 of No.~2).

Recall (§2, No.~5) that if \(\chi\) is continuous and \(\beta\) has support \(X\) , then \(\chi\) is a multiplier.

Let \(f\) be a locally \(\beta\) -integrable complex function on \(X\) , and let \(\mu\) be a measure on \(G\) . For every \(s \in G\) , the measure \(\boldsymbol{\gamma}(s)(f \cdot \beta)\) has base \(\beta\) since \(\beta\) is quasi-invariant. Therefore, if \(\mu\) and \(f \cdot \beta\) are convolvable, then \(\mu * (f \cdot \beta)\) has base \(\beta\) (§3, No.~2, Prop. 10).

DEFINITION 1. --- If \(\mu\) and \(f \cdot \beta\) are convolvable, \(\mu\) and \(f\) are said to be convolvable relative to \(\beta\) . Every density of \(\mu * (f \cdot \beta)\) with respect to \(\beta\) is called a convolution product of \(\mu\) and \(f\) relative to \(\beta\) and is denoted \(\mu * ^{\beta} f\) .

One omits \(\beta\) when no confusion is possible. Convolution for several measures on G and a function on X is defined in an analogous manner.

The various convolution products of \(\mu\) and f are equal locally \(\beta\) -almost everywhere. If \(\beta\) has support X and if there exists a convolution product of \(\mu\) and f that is continuous, then the latter is uniquely determined; it is then called the convolution product of \(\mu\) and f relative to \(\beta\) .

Let \(s \in G\) and let \(f\) be a locally \(\beta\) -integrable complex function on \(X\) . Then \(\varepsilon_s\) and \(f\) are convolvable, and

\[
\varepsilon_ {s} * (f \cdot \beta) = \boldsymbol {\gamma} (s) (f \cdot \beta) = (\boldsymbol {\gamma} (s) f) \cdot (\boldsymbol {\gamma} (s) \beta) = (\boldsymbol {\gamma} (s) f) \cdot \chi (s ^ {- 1}, \cdot) \cdot \beta ,
\]

therefore

\[
\left(\varepsilon_ {s} * f\right) (x) = \chi \left(s ^ {- 1}, x\right) f \left(s ^ {- 1} x\right) = \left(\gamma_ {\chi} (s) f\right) (x)\tag{2}
\]

locally \(\beta\) -almost everywhere.

Lemma 1. --- Let \(\mu\) be a measure on \(G\) . Then \(\chi\) is locally \((\mu \otimes \beta)\) -integrable, and the image of \(\mu \otimes \beta\) under the homeomorphism \((s, x) \mapsto (s, s^{-1}x)\) of \(G \times X\) onto \(G \times X\) is \(\chi \cdot (\mu \otimes \beta)\) .

We may suppose that \(\mu \geqslant 0\) . Let \(\mathbf{F} \in \mathcal{K}_{+}(\mathbf{G} \times \mathbf{X})\) . Then

\[
\begin{array}{l} \iint \mathrm{F} (s, s ^ {- 1} x) d \mu (s) d \beta (x) = \int d \mu (s) \int \mathrm{F} (s, s ^ {- 1} x) d \beta (x) \\ = \int d \mu (s) \int \mathrm{F} (s, x) d (\boldsymbol {\gamma} (s ^ {- 1}) \beta) (x) = \int d \mu (s) \int \mathrm{F} (s, x) \chi (s, x) d \beta (x). \end{array}
\]

Now, the function \((s,x)\mapsto\mathrm{F}(s,x)\chi(s,x)\) has compact support and is \((\mu\otimes\beta)\) -measurable. By Ch. V, §8, No.~3, Prop. 7, the preceding equality proves that this function is \((\mu\otimes\beta)\) -integrable and that

\[
\iint \mathrm{F} (s, s ^ {- 1} x) d \mu (s) d \beta (x) = \iint \mathrm{F} (s, x) \chi (s, x) d \mu (s) d \beta (x).
\]

This proves at the same time both assertions of Lemma 1.

PROPOSITION 1. --- Let \(\mu\) be a measure on G, \(f\) a locally \(\beta\) -integrable complex function on X. Suppose that the function \(s \mapsto f(s^{-1}x)\chi(s^{-1}, x)\) is essentially \(\mu\) -integrable except for a locally \(\beta\) -negligible set of values of \(x\) , and that the function \(x \mapsto \int |f(s^{-1}x)| \chi(s^{-1}, x) d|\mu|(s)\) , defined locally almost everywhere for \(\beta\) , is locally \(\beta\) -integrable. Then \(\mu\) and \(f\) are convolvable.

We may assume that \(f \geqslant 0\) and \(\mu \geqslant 0\) . Let \(h \in \mathcal{H}_{+}(X)\) . We are to prove that the function \((s, x) \mapsto h(sx)\) is essentially integrable for \(\mu \otimes (f \cdot \beta) = (1 \otimes f) \cdot (\mu \otimes \beta)\) (Ch. V, §8, No.~5, Prop. 10), that is, that \(\iint^{\bullet} h(sx)f(x)d\mu(s)d\beta(x) < +\infty\) (Ch. V, §5, No.~3, Prop. 3); it will clearly suffice to prove that there exists an \(a > 0\) such that for every compact subset \(K\) of \(G\) ,

\[
\iint^ {\bullet} h (s x) f (x) \varphi_ {K} (s) d \mu (s) d \beta (x) \leqslant a.
\]

By Lemma 1,

\[
\begin{array}{r l} \iint^ {\bullet} h (s x) f (x) \varphi_ {K} (s) d \mu (s) d \beta (x) \\ & = \iint^ {*} h (x) f (s ^ {- 1} x) \varphi_ {K} (s) \chi (s ^ {- 1}, x) d \mu (s) d \beta (x). \end{array}
\]

Now, the function \((s,x)\mapsto h(x)f(s^{-1}x)\varphi_{\mathrm{K}}(s)\chi(s^{-1},x)\) is \((\mu\otimes\beta)\) -measurable (Lemma 1) and has compact support. The preceding expression is therefore equal (Ch. V, §8, No.~3, Prop. 7) to

\[
\begin{array}{r l} & {\int^ {*} h (x) d \beta (x) \int^ {*} f (s ^ {- 1} x) \varphi_ {\mathrm{K}} (s) \chi (s ^ {- 1}, x) d \mu (s)} \\ & {\qquad \leqslant (\sup h) \int_ {\mathrm{S}} ^ {*} d \beta (x) \int^ {\bullet} f (s ^ {- 1} x) \chi (s ^ {- 1}, x) d \mu (s),} \end{array}
\]

where S denotes the support of h. Whence the proposition.

PROPOSITION 2. --- Let \(\mu\) be a measure on G, \(f\) a locally \(\beta\) -integrable complex function on X. Assume that one of the following conditions is satisfied:

\begin{enumerate}
\def\labelenumi{(\roman{enumi})}
\item
  \(f\) and \(\chi\) are continuous;
\item
  G operates properly in X and \(f\) is zero on the complement of a countable union of compact sets;
\item
  \(\mu\) is carried by a countable union of compact sets.
\end{enumerate}

If \(\mu\) and \(f\) are convolvable, then the function \(s \mapsto f(s^{-1}x)\chi(s^{-1}, x)\) is essentially \(\mu\) -integrable except for a locally \(\beta\) -negligible set of values of \(x\) , and one has, locally almost \(\beta\) -everywhere,

\[
(\mu * ^ {\beta} f) (x) = \int_ {G} f (s ^ {- 1} x) \chi (s ^ {- 1}, x) d \mu (s) = \int_ {G} (\gamma_ {\chi} (s) f) (x) d \mu (s).\tag{3}
\]

Let \(h \in \mathcal{K}(\mathrm{X})\) . Since \(\mu\) and \(f\) are convolvable, the function \((s, x) \mapsto h(sx)f(x)\) is essentially \((\mu \otimes \beta)\) -integrable. By Lemma 1, the function \((s, x) \mapsto h(x)f(s^{-1}x)\chi(s^{-1}, x)\) is essentially \((\mu \otimes \beta)\) -integrable. Under hypothesis (i) or (ii) of the statement, one then deduces that this function is \((\mu \otimes \beta)\) -integrable; for, in the first case it is continuous and one applies Prop. 3 of Ch. V, §1, No.~1, and in the second case it is zero outside a countable union of compact sets, and one applies Prop. 7, 2) of No.~2, loc. cit. By the Lebesgue--Fubini theorem,

\[
\begin{array}{c} \iint h (s x) d \mu (s) d (f \cdot \beta) (x) = \iint h (x) f (s ^ {- 1} x) \chi (s ^ {- 1}, x) d \mu (s) d \beta (x) \\ = \int h (x) d \beta (x) \int f (s ^ {- 1} x) \chi (s ^ {- 1}, x) d \mu (s), \end{array}
\]

the function \(x \mapsto g(x) = \int f(s^{-1}x)\chi (s^{-1},x)d\mu (s)\) being moreover locally \(\beta\) -integrable. One thus sees that

\[
\langle h, \mu * (f \cdot \beta) \rangle = \langle h, g \cdot \beta \rangle ,
\]

whence \(g = \mu *^{\beta} f\) .

Suppose now that \(\mu\) is carried by the union \(S\) of a sequence of compact sets. The function

\[
(s, x) \mapsto h (x) f (s ^ {- 1} x) \chi (s ^ {- 1}, x) \varphi_ {\mathrm{S}} (s)
\]

is essentially \((\mu\otimes\beta)\) -integrable, and zero outside a countable union of compact sets, hence \((\mu\otimes\beta)\) -integrable. Since \(\mu=\varphi_{S}\cdot\mu\) , the argument concludes as before.

Remark. --- The hypothesis (iii) of Prop. 2 is satisfied notably when \(\mu\) is bounded. For, for every \(n > 0\) , there then exists a compact subset \(\mathbf{K}_n\) of \(\mathbf{G}\) such that

\[
| \mu | (\mathrm{G} - \mathrm{K} _ {n}) \leqslant \frac {1}{n}
\]

(Ch. IV, §4, No.~7), and \(\mu\) is carried by the union of the \(K_{n}\) . More generally, let \(\rho\) be a lower semi-continuous finite function \(>0\) on \(G\) such that \(\rho(st) \leqslant \rho(s)\rho(t)\) ; if \(\mu \in \mathcal{M}^{\rho}\) , the hypothesis (iii) is satisfied; for, \(\rho \cdot \mu\) is bounded, and \(\mu\) is carried by the same subsets as \(\rho \cdot \mu\) since, on every compact subset of \(G\) , \(\rho\) is bounded below by a constant \(>0\) .

